% AM I 21, 29, 32; AF 06, 14, 15A; AM II 08–09.
\chapter{Ciągłość, różniczkowanie i całkowanie granic}
\label{ma2:ch:limits-calculus}

W rozdziałach \ref{ma2:ch:uniform-convergence}
i \ref{ma2:ch:function-series} ustaliliśmy warunki,
pod którymi granica funkcji ciągłych pozostaje
ciągła. Zbadamy teraz przemienność granicy
z innymi operacjami analizy. W przypadku
wyników udowodnionych w tomie I podajemy
odsyłacze zamiast ponownych dowodów.

\section{Zamiana kolejności granic}

\begin{twierdzenie}
\label{ma2:thm:interchange-limits}
Niech $X$ będzie przestrzenią metryczną,
$A\subseteq X$, $x_0$ punktem skupienia $A$,
a $E$ przestrzenią Banacha. Niech
$f_n:A\setminus\{x_0\}\to E$ zbiegają
jednostajnie do $f$ oraz niech dla każdego $n$
istnieje granica $a_n=\lim_{x\to x_0,\ x\in A\setminus\{x_0\}}f_n(x)$.
Wtedy istnieją obie granice
\[
 \lim_{x\to x_0,\ x\in A\setminus\{x_0\}} f(x)
 =\lim_{n\to\infty} a_n.
\]
\end{twierdzenie}
\begin{proof}
Z jednostajnej zbieżności wynika, że dla
$\varepsilon>0$ istnieje $N$ z
$\|f_n(x)-f_m(x)\|<\varepsilon/2$
dla wszystkich $m,n\ge N$ i wszystkich
$x\in A\setminus\{x_0\}$.
Dla ustalonych $n,m$ funkcja $f_n-f_m$ ma
granicę $a_n-a_m$ z ciągłości odejmowania,
a norma jest ciągła (twierdzenie
\ref{ma2:thm:metric}). Przechodząc do granicy
przy $x\to x_0$, otrzymujemy
$\|a_n-a_m\|\le\varepsilon/2$.
Z zupełności $E$ istnieje $a=\lim_n a_n$.
Dla $\varepsilon>0$ wybieramy $n$ tak duże,
aby $\|a_n-a\|<\varepsilon/3$ oraz
$\|f_n(x)-f(x)\|<\varepsilon/3$ dla wszystkich $x$.
Z definicji $a_n$ istnieje $\delta>0$,
dla którego $0<d(x,x_0)<\delta$, $x\in A$,
implikuje $\|f_n(x)-a_n\|<\varepsilon/3$.
Wtedy
\[
 \|f(x)-a\|\le
 \|f(x)-f_n(x)\|+\|f_n(x)-a_n\|+\|a_n-a\|
 <\varepsilon.
\]
\end{proof}
Wystarcza zbieżność jednostajna na pewnym
ustalonym przebitym otoczeniu $x_0$ w $A$.
Jeżeli granica $(a_n)$ już istnieje w $E$,
ostatni argument nie wymaga zupełności.
Bez jednostajności zamiana może zawodzić:
dla $f_n(x)=x^n$ na $[0,1)$ mamy
$\lim_n\lim_{x\to1^-} f_n(x)=1$ oraz
$\lim_{x\to1^-}\lim_n f_n(x)=0$ (por.
przykład \ref{ma2:ex:powers-pointwise}).

\section{Granica pod znakiem całki Riemanna}

Przypomnijmy pojęcie całki Riemanna z tomu I,
rozdział 29. Dla $a<b$ i podziału z punktami
pośrednimi $P=(a=t_0<\cdots<t_m=b;\
\xi_j\in[t_{j-1},t_j])$ piszemy
\[
 |P|=\max_j(t_j-t_{j-1}),\qquad
 R(f,P)=\sum_{j=1}^{m}(t_j-t_{j-1})f(\xi_j).
\]
Jeśli $f:[a,b]\to E$, gdzie $E$ jest przestrzenią
unormowaną, to całkowalność w sensie Riemanna
oznacza istnienie $I\in E$ takiego, że
dla każdego $\varepsilon>0$ istnieje $\delta>0$
i $\|R(f,P)-I\|<\varepsilon$ dla wszystkich
podziałów z punktami pośrednimi o $|P|<\delta$.
Piszemy $I=\int_a^b f(t)\,\mathrm dt$.
Granica jest jednoznaczna: jeżeli $I,J$
spełniają definicję, to dla $\varepsilon>0$
wybieramy podział o średnicy tak małej, aby
$\|R(f,P)-I\|<\varepsilon/2$ oraz
$\|R(f,P)-J\|<\varepsilon/2$. Wtedy
$\|I-J\|<\varepsilon$, a dowolność
$\varepsilon$ daje $I=J$.
Każda funkcja całkowalna w tym sensie jest
ograniczona. Istotnie, wybieramy jeden podział
o średnicy dostatecznie małej, aby wszystkie jego
sumy z dowolnymi punktami pośrednimi różniły się
od $I$ w normie o mniej niż $1$. Ustalamy punkty
pośrednie i zmieniamy tylko punkt w $j$-tym
podprzedziale. Różnica sum ma normę
$(t_j-t_{j-1})\|f(t)-f(\xi_j)\|<2$, więc
$\|f(t)\|\le\|f(\xi_j)\|+2/(t_j-t_{j-1})$
na tym podprzedziale. Maksimum skończonej
liczby takich ograniczeń daje ograniczoność.
Dla funkcji rzeczywistych jest to zatem definicja
równoważna definicji Darboux (AM I, twierdzenie
29.2.2, etykieta \texttt{ma1:thm:riemann-darboux-equivalence}).
W przypadku funkcji ciągłych o wartościach
w przestrzeni Banacha istnienie całki zostało
udowodnione w analizie funkcjonalnej, rozdział 15A,
twierdzenie o istnieniu całki przy zupełności
ciągowej (\texttt{thm:riemann-existence15a}).
Pełniejszą teorię całki wektorowej omówimy
w rozdziale 18.

\begin{twierdzenie}[Skalarna granica pod znakiem całki]
\label{ma1:thm:riemann-uniform-limit}
Niech $a<b$ i niech $f_n:[a,b]\to\R$ będą całkowalne
w sensie Riemanna i niech $f_n\to f$ jednostajnie.
Wówczas $f$ jest całkowalna w sensie Riemanna i
\[
 \left|\int_a^b f_n(t)\,\mathrm dt-
              \int_a^b f(t)\,\mathrm dt\right|
 \le (b-a)\sup_{t\in[a,b]}|f_n(t)-f(t)|
 \longrightarrow0.
\]
\end{twierdzenie}
\emph{Dowód:} AM I, twierdzenie 32.1.1, wraz
z oszacowaniem w jego dowodzie. Zachowano
oryginalną etykietę.

\begin{twierdzenie}[Granica całek wektorowych]
\label{ma2:thm:vector-riemann-uniform}
Niech $E$ będzie przestrzenią Banacha,
$a<b$, a $f_n:[a,b]\to E$ będą całkowalne
w sensie Riemanna. Jeżeli $f_n\to f$
jednostajnie, to $f$ jest całkowalna oraz
\[
 \left\|\int_a^b f_n(t)\,\mathrm dt-
                 \int_a^b f(t)\,\mathrm dt\right\|
 \le (b-a)\sup_{t\in[a,b]}\|f_n(t)-f(t)\|.
\]
W szczególności całki $f_n$ dążą do całki $f$.
\end{twierdzenie}
\begin{proof}
Oznaczmy $L=b-a>0$ i $I_n=\int_a^b f_n(t)\,\mathrm dt$.
Dla dowolnego podziału z punktami pośrednimi $P$
i funkcji $g,h$ mamy
\[
 \|R(g,P)-R(h,P)\|\le L\sup_{t\in[a,b]}\|g(t)-h(t)\|.
\]
Istotnie, różnica sum jest sumą wektorów
$(t_j-t_{j-1})(g(\xi_j)-h(\xi_j))$.
Z nierówności trójkąta jej norma nie przekracza
$\sum_j(t_j-t_{j-1})\|g(\xi_j)-h(\xi_j)\|$,
co po oszacowaniu każdej normy przez supremum
daje tezę, gdyż $\sum_j(t_j-t_{j-1})=L$.
Jeżeli $g,h$ są całkowalne, przechodzimy do
granicy przy $|P|\to0$, uzyskując tę samą
nierówność dla różnicy ich całek.
Dla $g=f_n$, $h=f_m$ z jednostajnego
warunku Cauchy'ego (twierdzenie
\ref{ma2:thm:uniform-cauchy}) wynika,
że $(I_n)$ jest ciągiem Cauchy'ego w $E$.
Z zupełności istnieje $I=\lim_n I_n$.
Ponadto $f$ jest ograniczona, bo z jednostajnej
zbieżności dla pewnego $n$ mamy
$\sup_t\|f(t)-f_n(t)\|\le1$, a $f_n$ jest ograniczona.
Ustalmy $\varepsilon>0$ i wybierzmy $n$ tak,
aby $\sup_t\|f_n(t)-f(t)\|<\varepsilon/(3L)$
i $\|I_n-I\|<\varepsilon/3$.
Z całkowalności $f_n$ istnieje $\delta>0$
takie, że dla $|P|<\delta$ mamy
$\|R(f_n,P)-I_n\|<\varepsilon/3$.
Wtedy
\[
 \|R(f,P)-I\|
 \le \|R(f,P)-R(f_n,P)\|
    +\|R(f_n,P)-I_n\|+\|I_n-I\|<\varepsilon.
\]
To jest definicja całkowalności $f$ i jej całki $I$.
Stosując początkowe oszacowanie do całkowalnych
$f_n,f$, otrzymujemy tezę.
\end{proof}
Liniowość całki w tej definicji wynika z liniowości
sum Riemanna i ciągłości działań w $E$:
dla całkowalnych $g,h$ oraz skalarów $\alpha,\beta$
mamy $R(\alpha g+\beta h,P)=\alpha R(g,P)+\beta R(h,P)$;
granica po obu stronach daje istnienie całki
kombinacji i jej wartość.

\begin{wniosek}[Całkowanie szeregu wyraz po wyrazie]
\label{ma2:cor:termwise-integral}
Niech $E$ będzie przestrzenią Banacha,
$a<b$, a $u_n:[a,b]\to E$ funkcjami całkowalnymi
w sensie Riemanna. Jeżeli $\sum_n u_n$
jest jednostajnie zbieżny do $S$, to $S$
jest całkowalna i
\[
 \int_a^b S(t)\,\mathrm dt
 =\sum_{n=1}^\infty\int_a^b u_n(t)\,\mathrm dt.
\]
Błąd po $N$ wyrazach ma normę nie większą niż
$(b-a)\sup_t\|S(t)-S_N(t)\|$.
\end{wniosek}
\begin{proof}
Sumy częściowe $S_N$ są całkowalne
z liniowości całki, a z założenia $S_N\to S$
jednostajnie. Zastosowanie twierdzenia
\ref{ma2:thm:vector-riemann-uniform} i ponowne
użycie liniowości daje równość oraz oszacowanie.
\end{proof}
Dla $a=b$ całka jest zerowa, a dla $b<a$
zmieniamy orientację; w oszacowaniach
używamy wtedy $|b-a|$.

\begin{przyklad}
\label{ma2:ex:termwise-geometric-integral}
Dla $0\le r<1$ szereg $\sum_{n=0}^\infty t^n$
jest jednostajnie zbieżny na $[0,r]$ do $1/(1-t)$
(przykład \ref{ma2:ex:function-geometric}).
Całkując wyraz po wyrazie, otrzymujemy
\[
 -\log(1-r)=\sum_{n=0}^\infty\frac{r^{n+1}}{n+1}.
\]
Dla $r>0$ użyto funkcji pierwotnych
$-\log(1-t)$ i $t^{n+1}/(n+1)$
(AM I, stwierdzenie 28.4.1)
oraz twierdzenia Newtona--Leibniza
(AM I, twierdzenie 29.6.3).
Dla $r=0$ obie strony równości są zerowe. Błąd po wyrazie
o indeksie $N$ nie przekracza
$r^{N+2}/(1-r)$, z oszacowania w przykładzie
\ref{ma2:ex:function-geometric} i wniosku
\ref{ma2:cor:termwise-integral}.
\end{przyklad}

\begin{przyklad}[Sama zbieżność punktowa nie wystarcza]
\label{ma1:ex:parameter-moving-spike}
Dla $n\ge2$ niech
\[
 f_n(t)=\begin{cases}
 n^2t,&0\le t\le1/n,\\
 2n-n^2t,&1/n\le t\le2/n,\\
 0,&2/n\le t\le1.
 \end{cases}
\]
Każda funkcja jest ciągła i $f_n\to0$ punktowo
na $[0,1]$, lecz $\int_0^1 f_n(t)\,\mathrm dt=1$.
Weryfikację rachunków i granicy zawiera AM I,
przykład 32.1.3. Zatem granica całek
nie jest całką granicy.
\end{przyklad}

\section{Granica pod znakiem pochodnej}

Pochodną $f:I\to E$ względem rzeczywistej
zmiennej definiujemy w normie $E$ jako granicę
ilorazów różnicowych. Na końcach domkniętego
przedziału rozumiemy ją jednostronnie,
o ile ją rozważamy.

\begin{lemat}[Oszacowanie przyrostu]
\label{ma2:lem:vector-increment}
Niech $E$ będzie przestrzenią unormowaną nad
$\R$ lub $\C$, $g:[u,v]\to E$ ciągła,
różniczkowalna w $(u,v)$, gdzie $u<v$, $M\ge0$.
Jeżeli $\|g'(t)\|\le M$ na $(u,v)$, to
\[
 \|g(v)-g(u)\|\le M(v-u).
\]
\end{lemat}
\begin{proof}
Dla $g(v)=g(u)$ lewa strona jest zerem,
a prawa jest nieujemna, więc teza zachodzi.
W przeciwnym przypadku traktujemy $E$ jako
rzeczywistą przestrzeń unormowaną i wybieramy
rzeczywisty ciągły funkcjonał liniowy $\ell$
o normie $1$, taki że
$\ell(g(v)-g(u))=\|g(v)-g(u)\|$.
Istnienie wynika z twierdzenia Hahna--Banacha
(analiza funkcjonalna, rozdział 06, wniosek
o funkcjonale normującym, \texttt{cor:norming}).
Skalarna funkcja $\ell\circ g$ jest ciągła
na odcinku i różniczkowalna wewnątrz, a jej
pochodna wynosi $\ell(g'(t))$ z ciągłości
$\ell$ i definicji pochodnej.
Z twierdzenia Lagrange'a (AM I, twierdzenie 21.3.1)
istnieje $c\in(u,v)$ z
\[
 \|g(v)-g(u)\|=\ell(g'(c))(v-u)
 \le\|g'(c)\|(v-u)\le M(v-u).
\]
Pierwsza nierówność wynika z $\|\ell\|=1$
i $v-u>0$.
\end{proof}
W przypadku skalarnym rzeczywistym jest to
szczególny przypadek twierdzenia 21.6.1 z AM I.
Nie zakładamy tutaj ciągłości pochodnej.

\begin{twierdzenie}[Jednostajna zbieżność pochodnych]
\label{ma2:thm:uniform-derivative}
Niech $I$ będzie przedziałem o niepustym wnętrzu,
a $E$ przestrzenią Banacha. Niech $f_n:I\to E$
będą ciągłe i różniczkowalne we wnętrzu $I$.
Załóżmy, że dla pewnego $x_0\in I$
ciąg $(f_n(x_0))$ jest zbieżny w $E$ i że
$f_n'\to g:\operatorname{int}I\to E$ jednostajnie
na $\operatorname{int}I$.
Wtedy $f_n$ zbiegają jednostajnie na każdym
ograniczonym podprzedziale $J\subseteq I$ do funkcji
$f:I\to E$ ciągłej i różniczkowalnej we wnętrzu,
a $f'=g$. Jeśli $I$ jest ograniczony, zbieżność
jest jednostajna na całym $I$.
\end{twierdzenie}
\begin{proof}
Z jednostajnej zbieżności pochodnych dla każdego
$\eta>0$ istnieje $N$ taki, że dla $n,m\ge N$
i $t\in\operatorname{int}I$ mamy
$\|f_n'(t)-f_m'(t)\|<\eta$.
Lemat \ref{ma2:lem:vector-increment} zastosowany
do $f_n-f_m$ na odcinku o końcach $x_0,x$
daje
\[
 \|f_n(x)-f_m(x)\|
 \le\|f_n(x_0)-f_m(x_0)\|+\eta|x-x_0|.
\]
Dla $x=x_0$ nierówność zachodzi bez lematu.
Na dowolnym niepustym ograniczonym
podprzedziale $J\subseteq I$ liczba $D=\sup_{x\in J}|x-x_0|$ jest skończona.
Dla $\varepsilon>0$ wybieramy
$\eta=\varepsilon/(2(1+D))$ oraz zwiększamy
$N$, aby pierwszy składnik był mniejszy niż
$\varepsilon/2$. Otrzymujemy jednostajny
warunek Cauchy'ego na $J$.
Z twierdzenia \ref{ma2:thm:uniform-cauchy}
istnieje granica jednostajna na $J$.
Granice na różnych podprzedziałach zgadzają się
w punktach wspólnych z jednoznaczności granicy
punktowej; definiują $f$ na całym $I$.
Każdy punkt $I$ ma ograniczone otoczenie względne,
więc ciągłość $f$ wynika z twierdzenia
\ref{ma2:thm:uniform-continuous-limit}.
Dla ograniczonego $I$ bierzemy $J=I$.
Otrzymujemy też oszacowanie błędu. Dla każdego
$n$ takiego, że
$A_n=\sup_{t\in\operatorname{int}I}\|f_n'(t)-g(t)\|$
jest skończone, zachodzi
\[
 \|f_n(x)-f(x)\|
 \le\|f_n(x_0)-f(x_0)\|+A_n|x-x_0|
 \qquad(x\in I).
\]
Istotnie, dla dostatecznie dużych $m$ norma
pochodnej $f_n-f_m$ jest nie większa niż
$A_n+A_m$, a $A_m\to0$ z jednostajnej
zbieżności pochodnych. Stosujemy oszacowanie
przyrostu do $f_n-f_m$ i przechodzimy po $m$
do granicy. W szczególności na ograniczonym
$J$ błąd jednostajny jest nie większy niż
$\|f_n(x_0)-f(x_0)\|+DA_n$.

Ustalmy $x\in\operatorname{int}I$ i $\varepsilon>0$.
Wybieramy $n$ tak, aby
$\sup_t\|f_m'(t)-f_n'(t)\|\le\varepsilon/3$
dla wszystkich $m\ge n$ i
$\|f_n'(x)-g(x)\|<\varepsilon/3$.
Dla $h\ne0$ z $x+h\in I$ lemat daje
\[
 \|(f_m-f_n)(x+h)-(f_m-f_n)(x)\|
 \le(\varepsilon/3)|h|.
\]
Przechodząc po $m$ do granicy, otrzymujemy
\[
 \left\|\frac{f(x+h)-f(x)}h-
             \frac{f_n(x+h)-f_n(x)}h\right\|
 \le\varepsilon/3.
\]
Przy ustalonym $n$ dla dostatecznie małego
$|h|>0$ iloraz różnicowy $f_n$ różni się
od $f_n'(x)$ w normie o mniej niż $\varepsilon/3$.
Nierówność trójkąta pokazuje, że iloraz
różnicowy $f$ różni się od $g(x)$ o mniej niż
$\varepsilon$. Stąd $f'(x)=g(x)$.
\end{proof}
Jeżeli wszystkie $f_n'$ są ciągłe, to $g$
jest ciągła jako granica jednostajna, zatem $f$
jest klasy $C^1$ we wnętrzu $I$.
Jeżeli $x_0\in\operatorname{int}I$, wystarczy
założyć zbieżność jednostajną pochodnych na
każdym zwartym podprzedziale zawartym we
wnętrzu $I$, aby otrzymać zbieżność $f_n$
jednostajną na zbiorach zwartych tego wnętrza
i równość $f'=g$ we wnętrzu. Istotnie, każdy
niepusty zwarty podzbiór wnętrza wraz z $x_0$
mieści się w zwartym podprzedziale wnętrza;
stosujemy na nim twierdzenie i łączymy
granice punktowe. W celu obliczenia pochodnej
w ustalonym punkcie wybieramy podprzedział
zawierający także jego otoczenie.
Nie wynika stąd zbieżność na końcach $I$,
jeśli do niego należą: na $[0,1]$ funkcje
$f_n(x)=x^n$ spełniają $f_n(0)=0$,
a na każdym zwartym podprzedziale $(0,1)$
ich pochodne $nx^{n-1}$ zbiegają jednostajnie
do zera. Dla $0<r<1$ mamy bowiem
$\sup_{0\le x\le r}|f_n'(x)|=nr^{n-1}\to0$
(iloraz kolejnych dodatnich wyrazów $nr^{n-1}$
dąży do $r<1$, więc szereg tych wyrazów jest
zbieżny z kryterium ilorazowego, AM I,
twierdzenie 9.6.1).
Mimo to granica punktowa na $[0,1]$ nie jest
ciągła w $1$. W tym przykładzie jako punkt
normalizacji można wybrać dowolny $x_0\in(0,1)$,
bo $x_0^n\to0$. Co więcej, dla $F_n(x)=nx^n$
i $x_0=1/2$ mamy $F_n(x_0)\to0$,
a na każdym zwartym podprzedziale $(0,1)$
pochodne $n^2x^{n-1}$ zbiegają jednostajnie
do zera: dla $0<r<1$ ciąg $n^2r^{n-1}$
dąży do zera z tego samego kryterium
ilorazowego (iloraz dąży do $r$).
Jednak $F_n(1)=n$ nie ma skończonej granicy.
Wariant lokalny nie zapewnia więc nawet
istnienia granicy na końcu przedziału.

\begin{wniosek}[Różniczkowanie szeregu wyraz po wyrazie]
\label{ma2:cor:termwise-derivative}
Niech $I$ będzie przedziałem o niepustym
wnętrzu, $E$ przestrzenią Banacha, a funkcje
$u_n:I\to E$ ciągłe i różniczkowalne
we wnętrzu $I$.
Jeżeli $\sum_n u_n(x_0)$ jest zbieżny dla pewnego
$x_0\in I$, a $\sum_n u_n'$ zbiega jednostajnie
we wnętrzu $I$, to $\sum_n u_n$ jest jednostajnie
zbieżny na każdym ograniczonym podprzedziale $J\subseteq I$
do różniczkowalnej we wnętrzu funkcji $S$ i
\[
 S'(x)=\sum_{n=1}^\infty u_n'(x)
 \qquad(x\in\operatorname{int}I).
\]
\end{wniosek}
\begin{proof}
Dla sum częściowych $S_N$ mamy
$S_N'=\sum_{n=1}^N u_n'$ z liniowości pochodnej.
Hipotezy są dokładnie hipotezami twierdzenia
\ref{ma2:thm:uniform-derivative} dla $(S_N)$.
\end{proof}

\section{Kontrprzykłady i zakres hipotez}

\begin{przyklad}[Jednostajna granica może nie być różniczkowalna]
\label{ma2:ex:uniform-limit-cusp}
Na $[-1,1]$ niech $f_n(x)=\sqrt{x^2+n^{-2}}$.
Funkcje są klasy $C^1$ oraz
\[
 0\le f_n(x)-|x|
 =\frac{n^{-2}}{\sqrt{x^2+n^{-2}}+|x|}
 \le\frac1n.
\]
W punkcie $x=0$ różnica wynosi $1/n$,
więc $\|f_n-|\cdot|\|_\infty=1/n$.
Zatem $f_n\to |x|$ jednostajnie. Funkcja $|x|$
nie jest różniczkowalna w zerze, ponieważ
ilorazy różnicowe z prawej i lewej strony
wynoszą odpowiednio $1$ i $-1$.
Pochodne $f_n'(x)=x/\sqrt{x^2+n^{-2}}$ zbiegają
punktowo do funkcji znaku z wartością $0$ w zerze.
Ta funkcja jest nieciągła w zerze, więc
zbieżność pochodnych nie jest jednostajna.
\end{przyklad}

\begin{figure}[htbp]
\centering
\begin{tikzpicture}[x=5cm,y=3.3cm,>=Stealth]
 \draw[->] (-1.08,0)--(1.1,0) node[right]{$x$};
 \draw[->] (0,-.06)--(0,1.28) node[above]{$y$};
 \draw[black,very thick] (-1,1)--(0,0)--(1,1);
 \foreach \n/\kol in {2/blue!75!black,4/teal!75!black,8/orange!85!black}{
  \draw[\kol,thick,domain=-1:1,samples=151]
      plot (\x,{sqrt(\x*\x+1/(\n*\n))});
  \fill[\kol] (0,{1/\n}) circle[radius=1.2pt];
 }
 % Legenda poniżej prawego ramienia granicy, bez przecinania krzywych.
 \draw[blue!75!black,thick] (.56,.40)--(.68,.40)
      node[right,black,font=\small] {$n=2$};
 \draw[teal!75!black,thick] (.56,.27)--(.68,.27)
      node[right,black,font=\small] {$n=4$};
 \draw[orange!85!black,thick] (.56,.14)--(.68,.14)
      node[right,black,font=\small] {$n=8$};
 \node[below] at (-1,0) {$-1$};
 \node[below] at (1,0) {$1$};
 \node[below left] at (0,0) {$0$};
\end{tikzpicture}
\caption{Gładkie funkcje z przykładu
\ref{ma2:ex:uniform-limit-cusp} i ich granica
jednostajna $|x|$ (czarna). Różnica w zerze
wynosi dokładnie $1/n$.}
\label{ma2:fig:uniform-limit-cusp}
\end{figure}
\FloatBarrier

\begin{przyklad}
\label{ma2:ex:derivative-mismatch}
Na $[-1,1]$ funkcje $f_n(x)=\sin(nx)/n$
zbiegają jednostajnie do zera, gdyż
$|f_n(x)|\le1/n$. Tymczasem $f_n'(0)=1$
dla każdego $n$, podczas gdy pochodna granicy
w zerze wynosi $0$. Nawet przy różniczkowalnej
granicy nie można bez dodatkowych założeń
zamienić pochodnej z granicą.
Co więcej, sama punktowa zbieżność pochodnych
także nie wystarcza. Na $[-1,1]$ funkcje
$h_n(x)=x/(1+nx^2)$ spełniają
$|h_n(x)|\le 1/(2\sqrt n)$, gdyż
$1+nx^2\ge 2\sqrt n|x|$.
Ich pochodne są równe
\[
 h_n'(x)=\frac{1-nx^2}{(1+nx^2)^2}.
\]
Dla $x=0$ mają wartość $1$, dla $x\ne0$
dążą do zera po podzieleniu licznika
i mianownika przez $n^2$.
Zatem $h_n\to0$ jednostajnie, pochodne
mają granicę punktową wszędzie, lecz
granica ta w zerze nie jest pochodną
funkcji granicznej.
\end{przyklad}

\begin{przyklad}[Znaczenie wartości w jednym punkcie]
\label{ma2:ex:derivative-anchor}
Na dowolnym niepustym przedziale funkcje stałe
$f_n(x)=n$ mają pochodne identycznie zerowe,
lecz nie mają granicy punktowej. Warunek
zbieżności wartości w pewnym punkcie
w twierdzeniu \ref{ma2:thm:uniform-derivative}
eliminuje tę dowolność stałej.
Na nieograniczonym przedziale $\R$ funkcje
$g_n(x)=x/n$ spełniają $g_n(0)=0$,
a pochodne $g_n'=1/n$ zbiegają jednostajnie
do zera. Funkcje $g_n$ zbiegają jednostajnie
na ograniczonych podprzedziałach, ale nie
na całej $\R$: dla $x=n$ mamy $g_n(n)=1$.
\end{przyklad}

W następnym rozdziale zastosujemy powyższe
twierdzenia do szeregów potęgowych.
